% a_evidence. Section file; new macros at the top; label prefix: evi
\section{Finite-order evidence}
\label{sec:evi}
\noindent\genusstatus{numerically supported} The finite-order checks recorded here are computations (exact arithmetic unless stated); they support the statements at the orders and genera given, and are not proofs. The code and outputs are in \texttt{code/} of this project.

\begingroup\sloppy

Nothing in this appendix is a proof step. No proof in this paper uses any of it. The arithmetic is exact, over $\bbQ$ or $\bbQ(\omega)$ or in a finite field where stated, unless stated otherwise. Each check compares two quantities computed in different ways; a \emph{control} is a deliberately altered version of the computation, and we say that it fails when it no longer agrees. All of this is at finite order: limits are collected in \cref{sec:evi-limits}.

Weights are written $(\alpha;\epsilon)=(\alpha_0,\alpha_1,\alpha_2;\epsilon_1,\epsilon_2)$ or, on $\TF$, $\epsilon=(\epsilon_1,\epsilon_2)$.

\subsection{Generic weights against localisation}
\label{sec:evi-gen}

\begin{itemize}
\item \emph{Genus-$2$ graph sum against localisation.} All $7$ stable graphs, rows from the canonical $R$-matrix of \cite{BS}, at $(\alpha;\epsilon)=(1,2,5;3,7)$, $(2,-3,7;1,4)$, $(1,3,8;5/2,-3/7)$, against localisation at $d=0,1$. Exact match in all three samples, for example $19189/100800$ and $34231/1680$ at the first. Poles in $z$ cancel through $y^2$. Symplectic defects vanish for $z$-order $\leq 5$, $y$-order $\leq 2$. The control $c_j\to-c_j$ on the constants $c_j$ with $\re^{-c_j}=h_{\wts_j}$ fails at both degrees in all samples. Near $\alpha=E/3$ the degree-$1$ localisation value tends to the refined BPS value at $d=1$: $(0,1)$ gives $48241/1680$ at $\epsilon=(3,7)$ and $81/80$ at $(1,1)$.
\item \emph{Limit at $d=2$, in floating point.} Near the diagonal ($\delta=10^{-2},10^{-3},10^{-4}$), the $y^2$ coefficient plus $6$ times the $y^1$ coefficient tends to the refined value at $d=2$: $-1596.42857$ at $(3,7)$, $-54$ at $(1,1)$, $53.16494$ at $(5/2,-3/7)$.
\item \emph{Three-dimensional stationary phase.} Through $z^3$, $y^2$, for two weight sets and three colours, the formal expansion agrees with the ratio of prefactors times $\prod_3\Theta(-z/a)$ exactly in all $6$ cases. The control with $\Theta(+z/a)$ fails in all $6$.
\item \emph{Value at $d=0$.} The value $\wprod_i(p_3-p_1^3)/17280$ reproduces $19189/100800$, $158713/290304$, $-10457963/29030400$.
\item \emph{Rows derived without the functions $D$, $\Ey$.} The structure constant of the second row was solved from $z\to\infty$ of $\Fsol(H)|_{P_j}$, through $y^3$; $A^2D=X$ and $\Ey=(D-X)/2$ through $y^3$, one weight set. Graph sum against an independently written localisation, exact: $g=2$, $d\leq 3$ at $(1,2,5;3,7)$, $(2,-3,8;1,4)$; $g=3$, $d\leq 2$ at $(1,2,5;3,7)$; $g=4$, $d\leq 1$ at $(1,2,5;3,7)$, which reaches $z^7$, that is, the constant $a_4$. With legs: $\langle H^a\rangle_1,\langle H^a\rangle_2$ through $d=2$ at three weight sets. The Hodge-vertex form equals the point form at $g=2,3$ through $y^2$. The Wick expansion $\Phi(\alpha,Y;z)$ equals $(\Delta^{1/2}\Rt_{0i})\,h_{(3\alpha_i,\tau,\tau')}$ through $z^3$, $y^2$, three colours, two weight sets; the control $h^{-1}$ fails. Constants: $\GW_{2,0}^{\TF}=E^2/640$ and $\GW_{3,0}^{\TF}=E^4/161280$ by exact interpolation at $\epsilon=(3,7),(1,1),(5/2,-3/7)$.
\item \emph{Rows from the $I$-function.} Rows from the $I$-function at $P_j$, the row recursion and the normalisation, as exact $y$-series. Poles in $z$ cancel; $2\Ey=D-X$. Graph sum minus localisation is $0$ at $y^0,\dots,y^2$ ($g=2$, two $\alpha$) and $y^0,\dots,y^3$ ($g=3$). Controls (sign flip of $c_j$, dropping the weights $-\epsilon_1,-\epsilon_2$, no $h$) fail, already at $y^1$. At $g=3$, $y^0$ the sign flip is invisible, because $\dim\Mbar_3$ is even.
\item \emph{The defect of the flatness proof.} The defect $\Mdef$ of \cref{prop:gen-flat} was computed directly from the rows of \cref{sec:gen}, with the structure constants solved from the expansion of the $I$-function at $z=\infty$. One has $\Mdef=\mathsf{D}\Mdef_{\mathrm{rat}}\mathsf{D}^{-1}$ with $\mathsf{D}=\mathrm{diag}(\wprod_j^{1/2})$ constant. $\Mdef_{\mathrm{rat}}=0$ over $\bbQ$ at $6$ rational samples through $z^5$, $y^3$, and through $z^7$, $y^5$ at $(1,2,5;3,7)$. $\Mdef=0$ in $\mathbb{F}_p$, with the square roots as elements of $\mathbb{F}_p$, at $5$ samples, $p\approx 2^{89}$, and through $z^9$, $y^6$ at $(1,2,5;3,7)$, $p\approx 2^{127}$. Controls $\Rm^*$, $\Rm(-z)$, $H\qprod\to-H\qprod$, $\theta\to-\theta$ fail, $\Rm^*$ first at $\psi^1$, $y^1$. The replacements $h\to h^{-1}$ and no $h$ give $\Mdef=0$, as \cref{lem:gen-ambig} says.
\item \emph{Class level.} Graber--Pandharipande localisation with $\psi$ classes pulled back from $\Mbar_{g,n}$ against the graph sum of $\Rm.\TFT$ with legs, exact over $\bbQ$: $150$ comparisons per weight set at $(g,n)=(1,1),(1,2),(2,1),(2,2)$ through $d=3$ at $(1,2,5;3,7)$, $(1,3,8;5/2,-3/7)$, $(3,7,-11/2;5/3,-2/7)$; $n=3$ at $g=1,2$ through $d=2$ at the first set; $g=3$, $n=1$ through $d=2$ at the first set and $d=1$ at the second; no mismatch. Also through $d=2$ at two weight sets, no mismatch: $\lambda_1$ insertions; Hodge-vertex form with legs; descendants at $k\leq 6$ from ancestors through the genus-$0$ fundamental solution. Divisor equation with the bubble term on both sides, $n=1\to2$ through $d=3$, $n=2\to3$ through $d=2$, two weight sets, no failure; the same equation without the bubble term fails. Controls that fail: legs at $+\psi$, $\Rm^*$, $\Rm(-z)$, $h^{-1}$, no $h$.
\item \emph{Comparison with \cite{LP18}.} $R_1=(1-L^2)/18$, $R_2=(1-24L-2L^2+25L^4)/648$ compared exactly with the $I$-function at $H=\lambda_0=1$ through $y^8$; a control fails. Here $L=(1+27y)^{-1/3}$ and $\lambda_0$ are in the notation of \cite{LP18}, not of this paper.
\item \emph{Mirror theorem and differential equation.} Direct graph sum of virtual localisation on $\Mbar_{0,n}(Z,d)$, $n=1,2$, $d\leq 4$, independent of the $I$-function, at $\alpha=(2/7,5/11,-3/13)$, $(1/3,-4/5,11/2)$, and $(1/7,3/5,11/13)$ with $e_1(\alpha)=\epsilon_1+\epsilon_2$. Results at all three: the mirror theorem at $P_0,P_1,P_2$ through $y^4$ at a rational $z$; the $1/z$ coefficients of the expansion of the $I$-function; the pairing; $\Fsol(1)=\Jsm/z$ and $\Fsol(H)=\thQ\Jsm$ from two-point sums; the equation $z\thQ\Fsol(H^a)=\sum_cC^c_{1a}\Fsol(H^c)$, $a=0,1,2$, through $\cQ^4$, with $C$ from three-point invariants. Control: dropping the mirror map fails. Values $\GW_{0,d}\epsilon_1\epsilon_2=3,-45/8,244/9,-12333/64$ at every point, and $\langle H^2\rangle_{0,1,d}\,\epsilon_1\epsilon_2/e_1(\alpha)=-5,87/4,-1544/9,27671/16$.
\end{itemize}

\emph{Sampling warning.} The sample $\alpha=(2,-3,7)$ has $w_{\mathsf{F}}+w_{\mathsf{F}'}=0$ on fixed graphs with $d\geq 2$, because $\alpha_1-\alpha_0=-(\alpha_2-\alpha_0)$, so term-by-term localisation is singular there. It was used only for $d\leq 1$. It should be avoided at $d\geq 2$.

\subsection{Rows, traces and the specialisation to \texorpdfstring{$\TF$}{TF}}
\label{sec:evi-spe}

\begin{itemize}
\item \emph{Closed form against Wick expansion.} $\Phi(\alpha^*,Y;z)$ by three-dimensional Wick expansion equals $\Theta(-z/E)\Srow(v,z)$ through $z^3$ at $(\epsilon;y)=(3,7;1/5)$, $(1,1;-2/7)$, $(5/2,-3/7;3/11)$; the control $\Theta(+z/E)$ fails. With a separately written implementation: exactly through $z^9$ at $(\epsilon;Y)=(3,7;-2/3)$, $(2,-5/3;5/7)$. And $\Phi(\alpha,-\alpha_i;z)=\prod\Theta(-z/a)$ through $z^4$ at three colours, $\alpha=(1,2,5)$.
\item \emph{Genus-$2$ total as a trace.} The total as $\Tr\otimes\Tr$ in $\ringC\otimes\ringC$ at $\alpha^*$, with rows only from $\Phi$ and no $y$-series, equals the genus-$2$ series of \cite{BS} exactly at the same points and $\sfF\in\{0,1,-2/3,5/3,11/7\}$. The control $\Theta(-z/\epsilon)$ fails at the one value tested per sample. The formula \cref{eq:spe-pair} equals the pairing from explicit rows at $A=7/4$, $M,b\leq 3$.
\item \emph{Behaviour near $\TF$ at $g=2$.} At fixed $y$, $\sfF$, the total at $\alpha^*+\delta(1,2,-3)$ minus the value on $\TF$ is $1.02\cdot 10^{-3}$, $1.03\cdot 10^{-5}$, $1.03\cdot 10^{-7}$, $1.03\cdot 10^{-9}$ for $\delta=10^{-1},\dots,10^{-4}$: no blow-up at fixed $y$. At generic $\alpha$, $y_0=10^{-25}$, the trace total minus $\GW_{2,0}$, over $y_0$, equals $\GW_{2,1}$ up to $\cO(y_0)$; the control with wrong $h$ gives $0.3$--$0.45$. Along $\alpha=\alpha^*+s(1,2/3,-5/7)$, $s=1,\dots,13$: fifth differences of $\epsilon_1\epsilon_2\GW_{2,1}$ and eighth differences of $\epsilon_1\epsilon_2e_3(\alpha)\GW_{2,0}$ vanish; the control without $e_3(\alpha)$ fails. Values on $\TF$ at $s=0$ by interpolation: $\GW_{2,1}=48241/1680$, $81/80$, $-13093109/26342400$; $\GW_{2,0}=5/32$, $1/160$, $841/125440$. The $y$-series of the genus-$2$ series of \cite{BS} has $y^0$ coefficient $E^2/640$, $y^1$ the refined value at $d=1$, and $y^2+6y^1$ the refined value at $d=2$. Further $\mathrm{N}(W)|_{y=0}=-1440$ at $\alpha=(1,2,5)$, and $\mathrm{N}(W)=E^9y^2/\xi^2$ at $y=1/5,-2/7$, $E=10$, with $\mathrm{N}(W)$ the norm from $\ringC$.
\item \emph{$\Phi$ against the $I$-function.} At $\alpha=(1,2,5)$, $\Phi(\alpha,Y;z)/\Phi(\alpha,-\alpha_0;z)$ against the stripped $I$-function, for $z$-order $\leq 4$ (and $\leq 5$, three colours, $K=3$), truncation $K=2,\dots,5$ at $y=1.5\cdot 10^{-8}$: differences fall with $K$ as truncation only. With a separate implementation, through $z^4$: $\Theta(z/\epsilon_1)\Theta(z/\epsilon_2)\Phi$ at the root $Y_j(y)$ equals the $I$-function row times $h_{\wts_j}$ up to $\cO(y^6)$, scaling confirmed at $y=10^{-5},10^{-7}$.
\item \emph{Totals on $\TF$.} At $g=2$ the total equals the genus-$2$ series of \cite{BS} exactly at $(\epsilon;t;\sfF)=(3,7;1/2;1)$, $(1,1;-2/3;5/3)$, $(5/2,-3/7;3/11;-2/3)$. At $g=3$, $\Tsumb_3(F(y))$ at $t=10^{-3},10^{-4},10^{-5},10^{-7}$ agrees with $\GW_{3,0}+\cQ\GW_{3,1}+\cQ^2\GW_{3,2}+\cQ^3\GW_{3,3}$, with refined BPS numbers as external control, with stable residual, the next coefficient. $\GW_{3,0}=125/2016$ at $(3,7)$, $707281/6195732480$ at $(5/2,-3/7)$, $1/10080$ at $(1,1)$. Flipping $\Theta(-z/E)$ or $\Theta(z/\epsilon)$ fails. $\Hg_{\mathsf{v}}=h_{\mathsf{v}}.1$ exactly at $g=2,3$, for three- and five-element weight sets. $\GW_{g,0}^{\Tmax}=\sum_i\wprod_i^{g-1}\int\Hg_{\wts_i}$ at $g=2,3$: polynomiality along a line and interpolation equal the formula of \cref{prop:spe-const} at $(3,7)$, $(5/2,-3/7)$.
\item \emph{Direct localisation against the point-frame sum.} Direct virtual localisation of $Z$, with no mirror symmetry and no $R$-matrix, at generic $\alpha$; exact polynomiality in $\alpha$; interpolation to $\alpha=E/3$. $\Tsumb_g[F(y)]$ at $t=\pm1/50,\pm1/500$: the quantity $(\Tsumb_g-\sum_{d<\mathsf{d}}\cQ^d\GW_{g,d})/\cQ^{\mathsf{d}}$ tends from both signs of $t$ to $\GW_{g,\mathsf{d}}$: $g=2,3$, $\mathsf{d}=3$ at $(3,7)$, $(1,2)$, and at $(5/2,-3/7)$ for $g=2$. For example, at $g=3$, $(3,7)$: $2827874.165$, $2827876.721$ at $t=\pm1/500$ against $2827875.443$. At $g=4$, $\mathsf{d}=2$, $(3,7)$, $t=1/500$: $-289820.888$ against $-289821.616$. Controls (rows $\hb^{-1}$, rows without $\hb$, $F+1/9$) fail.
\end{itemize}

\subsection{Frame change, gap and \texorpdfstring{$\tcf$}{tcf}}
\label{sec:evi-fra}

\begin{itemize}
\item \emph{Engine controls for Hodge integrals.} From the Witten--Kontsevich recursion: $1/1152$, $1/82944$, $7/240$. Mumford's recursion: $\lambda_1^3=1/2880$ on $\Mbar_2$, $\lambda_1^6=1/90720$ on $\Mbar_3$, $3!\,\mathrm{ch}_3=-1/5760$ on $\Mbar_2$. Graph counts $7$, $42$, $379$ at $g=2,3,4$.
\item \emph{Frame change.} At fixed $t$ and numeric $\sfF$, rows from the closed form, at $(\epsilon;t)=(3,7;1/5)$, $(1,1;-2/7)$, $(5/2,-3/7;3/11)$, $\sfF\in\{0,1,-2/3,5/3,11/7,-7/2,2/9\}$. $g=2$: in $21$ of $21$ cases $\Tsumb_2=\Gsum_2$, equal to previously computed tables and to the genus-$2$ series of \cite{BS}. $g=3$: in $21$ of $21$ cases $\Tsumb_3=\Gsum_3$, equal to previously computed tables. Since $\deg_{\sfF}\leq 6$, this is the polynomial identity at those points. All graph values differ between frames ($7$ of $7$, $42$ of $42$), so the identity is for the total only. Controls (Hodge vertices with rows $\hb\Ph$; point vertices with rows $\hb^{-1}\Ph$) fail in every case.
\item \emph{Frame change at $g=4$.} $379$ graphs: point frame equals Hodge frame equals a previously computed table at $(3,7;t=1/5;\sfF=1)$ and $(1,1;t=-2/7;\sfF=-2/3)$. The value at the first is $73603363895647521165850828027/67984534801612800000000$.
\item \emph{Gap as Laurent series.} Tables of the genus $2,3,4$ totals at $\Fc=\Fcf$ as exact Laurent series in $\xi$, against $(-3)^{g-1}\beta_{2g-2}/\tc^{2g-2}$, at $\epsilon=(3,7)$, $(1,1)$, $(5/2,-3/7)$, $(1/3,1)$, $(4/5,1)$: every polar coefficient of the difference is $0$ at $g=2,3,4$. Leading values reproduce $-169/27216$ at $g=3$, $\mathsf{x}=1/3$, and $-1414951/224000000$ at $g=3$, $\mathsf{x}=4/5$. The control $\Fc\to\Fc+1/9$ fails in all cases; $\Fc\to\Fc+\xi/5$ fails at $g=3,4$ and does not fail at $g=2$ (\cref{sec:evi-limits}). First regular value, $g=2$ at $(3,7)$: $-988747/4898880$.
\item \emph{Graph sum for $\beta$.} The graph sum of \cref{rmk:reg-beta} equals $\ZZ[g,0]=\beta_{2g-2}/(-\epsilon_1\epsilon_2E^3)^{g-1}$ exactly at $g=2,3,4$, $\epsilon=(3,7),(1,1),(5/2,-3/7)$; with $\hb(+z)$ the sum differs ($g=2,3$). The limit coefficients of $\beta$ at $E=0$ and $\epsilon_1\to0$ equal those of \cite{BS} at $g=2,3,4$: $-1/80$, $1/112$, $-3/160$; $-7/1920$, $-31/17920$, $-381/143360$. Also, at $g=2$, $(\epsilon_1,\epsilon_2)=(7/3,5/2)$, all $7$ graphs, the sum of \cite{BS} is $-152041/4097352000=\ZZ[2,0]$, and this times $-E^3\epsilon_1\epsilon_2$ is $\beta_2=152041/6220800$.
\item \emph{Frame change, independent graph enumeration.} $\Tsumb_2=\Gsum_2$, equal to the genus-$2$ series of \cite{BS}, at random $(\epsilon,t,\sfF,\cone)$ in $\bbQ(\omega)$; independence of $\cone$ confirmed. $\Tsumb_3=\Gsum_3$ at $t=\pm1/50,\pm1/500$ with $F(y)$. Symplecticity with free $\cone,\sfF$ exact through $z^8$ on $\TF$. Controls fail.
\item \emph{Gap by direct evaluation.} Exact evaluation of $\Gsum_g[\Fcf]$ at rational $t=-1/3+\delta$, with no previously computed data. $\Gsum_g-(-3)^{g-1}\beta_{2g-2}/\tc^{2g-2}$ stays bounded and converges: $g=2$ ($\delta=10^{-2},\dots,10^{-5}$), $g=3$ ($10^{-2},\dots,10^{-4}$) at three pairs $\epsilon$, and $g=4$ at $(3,7)$: $0.3303$, $0.3858$ at $\delta=10^{-2},10^{-3}$. Per colouring ($g=2,3$): the maximum of $|\GCont(\Gamma,\col)|$ over graphs with an edge is bounded; the colour-$1$ vertex is bounded; the colour-$0$ vertex minus its polar term is bounded. Controls: $F+1/9$ diverges; $F+\xi/5$ diverges at $g=3$; $\xi$ in place of $\tc$ diverges.
\item \emph{The coordinate $\tcf$.} $\tc$ minus the expression of \cite{BS} on $\log y=\log(-y)+\ri\pi$: $2\cdot 10^{-16}$, $-7\cdot 10^{-16}$, $-4\cdot 10^{-16}$ at $y=-1/54,-1/81,-1/270$ (double precision); below $7\cdot 10^{-16}$ also at $y=-1/40$. With $60$ digits: agreement to $10^{-58}$ at $y=-1/54,-1/81,-1/40$ ($10^{-38}$ at $-1/270$, truncation), and $\Acf+\wcf=10^{-56}$. Other branch at $y=-1/54$: real part $-4.6936$, imaginary part $6.6811$, consistent with the identity in \cref{lem:fra-tcf}.
\item \emph{Barnes's second lemma.} Parameters $(2/3,2/3,2/3,4/3,5/3)$, trapezoid rule on $\mathrm{Re}\,\sigma=-1/2$, step $1/120$, $60$ digits: left side minus $(2\pi/\sqrt3)^3$ is $8.1\cdot 10^{-47}$; generic control $(0.3,0.45,0.7,0.2)$: $8.6\cdot 10^{-32}$; the wrong value $e'=2$ fails ($47.98$ against $67.04$).
\end{itemize}

\subsection{Rows, leaves and the vertex identity}
\label{sec:evi-rows}

\emph{Caveat.} By the scaling law of \cref{rmk:lea-scale}, row and edge factors at different pairs $\epsilon$ differ only by $\epsilon_1\epsilon_2$ and the scale $E$. So row and edge checks ``at $n$ weight pairs'' below are one test of the scale-free object plus a test of the scaling; they do not probe non-special framing. Checks of quantities on the Hodge side do depend on $\tau$.

Exact rational arithmetic unless stated.
\begin{itemize}
\item Form of row $1$ through $\tc^7$, $z^9$ at $4$ weight pairs. Nine edges regular; rank one at $\tc=0$; support of loops and legs for $m+n\leq 8$, $\tc^4$; the value $\Fc(0)=-2/9$ is needed.
\item Per colouring: genus $2$ ($6$ graphs with an edge, $2$ weights) and genus $3$ ($41$ graphs with an edge, $2$ weights), all regular colouring by colouring; in the single-vertex graph only the singular colouring has a pole, which equals the value of \cref{cor:reg-sv}.
\item Leg and loop configurations have no polar term at $g=2,3$; $42$ heavier genus-$3$ configurations were skipped, since they need rows beyond $z^9$.
\item Frobenius cross-check of the rows: only simple poles at $z=-E/j$.
\item Previously computed rows equal the closed form at finite rational values of $Y$, hence at finite $\xi$, $z^2,\dots,z^9$, $4$ weight pairs; the control $z\to-z$ fails.
\item The first-order system of the rows at $\cone=\ccf$ holds at the coefficients of $z^1,\dots,z^9$ through $\tc^4$; away from $\cone=\ccf$ it fails at $z^2$.
\item No term linear in $\psi$ in $\brow_2$ through $\tc^7$.
\item \Cref{lem:lea-res} reproduces $\Edgeh_{00}$ exactly: $45$ entries, $\tc^{-2},\dots,\tc^{6}$, $3$ weight pairs; controls away from $\cone=\ccf$ fail.
\item $\Ey=(D-X)/2$ to $y^{40}$.
\item Previously computed rows satisfy the Picard--Fuchs equation exactly through $z^{10}$ at $6$ points. This cannot detect a wrong factor depending only on $z$. The closed form satisfies it exactly through $z^{19}$.
\item \Cref{prop:lea-H2} and \cref{thm:lea-M} at finite orders; consistent section $\cone=-\wcf(1+\mathsf{a}\xi)$ at $\epsilon=(7,3)$, $\mathsf{a}=1/10$ and $(4,9)$, $\mathsf{a}=-3/7$ ($k=-1,\dots,5$). Direct rank-one check of $\Edgeh_{0j}$ at $\tc=0$ in all entries. Recurrence for $r_l$ against Wick numbers; the Picard--Fuchs operator identity.
\item \Cref{lem:reg-string,lem:reg-L0} on the trivial cohomological field theory at genus $0$ at $E=7/3$, and at genus $0$ and $1$ by a separate recursion from $\langle\tau_1\rangle_1=1/24$, in $11$ cases. These are controls of the algebra, not of the Hodge statement. The controls $J/E$ instead of $J/E^2$, and $\dg$ instead of $\dg+2$ at $h=1$, fail.
\item Single-vertex pole against previously computed values, $5$ of $5$ cases: $g=2$ at $\mathsf{x}=1/3$: $-83/6480$; $\mathsf{x}=1/8$: $-9773/327680$; $g=3$ at $\mathsf{x}=1/3$: $-169/27216$; $\mathsf{x}=1/8$: $-1666067/117440512$; $\mathsf{x}=4/5$: $-1414951/224000000$.
\item $\Ltwo(\theta\Sinf)=0$ through $V^{14}$; a perturbed $r_2$ gives a non-zero result.
\item $(-3)^{g-1}\beta_{2g-2}(\mathsf{x},1)$ against previously computed local $\bbP^2$ values at $g=2,3,4$: agreement. These values are total leading polar coefficients. At $g=4$ the per-graph line for the single-vertex graph shows the gap followed by zeros.
\item $\Ltwo A=0$, $A^2D=X$, $2\Ey=D-X$ and the limit relations through $y^{60}$; the expansion of the $I$-function reproduced from its definition for $d\leq 6$ at a generic rational $\alpha$. Also, from the $I$-function expanded directly at $4$ values of $\alpha$, including $\TF$: $A^2D=X$, $2\Ey=D-X$ through $y^{22}$. Control fails.
\item The pairing of rows equals $\delta_{ij}$ on the closed form alone, all $i\leq j$, through $z^6$ at $(1,1)$, $(5,2)$, and through $z^5$ at $(7,-3)$, $(1/2,3/5)$. Controls ($r_2$ perturbed, $\eta^{00}$ doubled) fail. In $\mathbb{F}_p$, $p=2^{61}-1$, random section $\cone$, random $t_0$: for all colour pairs through $z^{14}$ at $9$ weight pairs. Controls ($\dtwo$ doubled, $r_3$ perturbed, $\eta^{00}$ doubled) fail.
\item $\dtwo=e_1(\alpha)\Ey$ and $\xi\cone^2=1/D$ at $\cone=A$ through $t^{24}$ (Laurent truncation); \cref{def:set-rows} equals the row recursion of \cite{BS} through $z^4$ up to $t^4$ at colour $0$ and $\epsilon=(5,2),(1,1)$.
\item In $\mathbb{F}_p$, at generic $\alpha$, with $\eta_{ab}$ as in \cref{prop:gen-qde} and an independent WKB solution of the Picard--Fuchs equation with rows from the row recursion: the pairing of different colours vanishes and $\theta$ of the pairing at equal colours vanishes, through $z^5$. Controls fail.
\item \emph{Vertex identity.} \Cref{eq:val-star} on the Schur side, $p_1^n$, $n\leq 7$ at $\tau=2/7,-5/3$; $\ZZ_h$ extracted from Mari\~no--Vafa data, $h=2,\dots,6$, equal to $\beta_{2h-2}$; \cref{eq:val-E1} with dual numbers in $\tau$. In monomials $x^a\pw^b$, $x$-order $\leq 12$: grade $\leq 1$ of \cref{eq:val-star} vanishes with $\ZZ_h=\beta_{2h-2}/\st^{h-1}$. Perturbing $\ZZ_3$ gives non-zero grade $1$.
\item Values free of the Mari\~no--Vafa formula, from Mumford's and Faber's recursions and the Witten--Kontsevich recursion: $\st^{h-1}\ZZ_h=\beta_{2h-2}(-\tau,1+\tau)$ for $h=2,3,4$ at $\tau=3/11,-7/4,5/2$, and for $h=5$ at $\tau=3/11$. The control $\beta\cdot\st$ fails.
\item Slice normalisation against Faber integrals for $(h,a)$ up to $(3,2)$; \cref{eq:val-E1} with no mismatch for $N^1$ down to $N^{-33}$ at $\tau=3/11,-7/4,13/5$. Kernel control: adding $(5/3)f_*((1+\tau)\mathsf{t})^4$, which breaks the symmetry in $\tau$, still solves \cref{eq:val-E1}, so the symmetry step is genuinely needed.
\item Modulo $2^{61}-1$: \cref{eq:val-star} on $p_1^n$, $n\leq 12$; $\ZZ_h$ for $h=2,\dots,8$ at $\tau=3/11,-7/4,13/5$ by a polynomial fit in $\pw$, equal to $\beta$.
\item With the next grade and the higher grades as arbitrary random data, grade $0$ of \cref{eq:val-star} vanishes and grade $1$ equals the operator of \cref{eq:val-E1} applied to the data minus its right side, exactly. Controls: perturbing the logarithmic coefficient or the genus-$1$ constant fails; perturbing the coefficient $-1/4$ of the genus-$0$ term changes nothing, as the kernel of that operator predicts.
\end{itemize}


\emph{Data on the vanishing, earlier computations.} Exact values of $\Mirr(3)$, $\Mirr(4)$; all $194$ points at genus $7$ are zero, with $77$ controls non-zero; a comparison of $\beta$ with values of the gap in $44$ cases through $h=12$.

\subsection{The vanishing theorem}
\label{sec:evi-van}

All exact unless stated.
\begin{itemize}
\item \emph{Formulas of \cref{sec:van}, fresh evaluation:} the Schur side of \cref{prop:van-real}, its Laguerre form, \cref{eq:van-Mk}, the target coefficient, at $\tau=2/7,-5/3$, $G=2,3$, $k_0\leq 4$, $n=2G,2G+1$, $48$ relations each: no failure. Target response exact, $24$ of $24$. Flipping a sign, or $\gamma_{\tau}\to-\gamma_{\tau}$, breaks $13$ of $24$; the rest are terms that vanish anyway. $\Xit_2(3)=1367/186624$, $\Xit_3(5)=205027/47029248$ at $\tau=2/7$.
\item \emph{\Cref{eq:van-real} against an independent Hodge integral engine:} from the Witten--Kontsevich and Mumford recursions, implemented anew. First run: $g\leq 4$, $150$ pairs $(g,\mu)$, $1387$ comparisons. Second run: $g\leq 5$ at $\ell(\mu)\leq 2$, $g\leq 4$ at $\ell(\mu)=3$, $595$ comparisons, several $\tau$ including zeros of $\Pi$. No mismatch; nine kinds of control all fail.
\item \emph{The formula of \cite{LLZ03}, eqs.~(1)--(3), literally in $\lambda_{\mathrm{MV}}$ over $\bbQ(\ri)$:} $90$ cases match; an incorrectly transcribed variant of the formula fails $66$ of $90$.
\item \emph{All-profile matching directly from the Schur sum:} $g\leq 3$, and $g\leq 6$ for the empty profile; no mismatch. Controls fail.
\item \emph{\Cref{lem:van-dil} and \cref{lem:van-ops}:} genus $0$ to $4$, non-integer leaves; no failure.
\item \emph{The finite formula against the definition,} $h=2,3,4$; zeros on the Schur side at genus $4$ to $7$; $\Xit_4(7)=193375/214990848$ at $\tau=2/7$ is non-zero, as expected out of range. No difference.
\item \emph{\Cref{eq:van-bil}} for $k=2,\dots,7$ on generic content-product $\tau$-functions, and for $k\leq 10$ at low weight, degree $8,9$: all vanish; controls (swapped order, $\tau_{-1}$, a wrong coefficient) are non-zero. On $\MVZ$ as Laurent series in $x$: $\tau=2/7,-5/3$, $k\leq 6$, weight $\leq 8$, order $30$: $114$ coefficients, no failure.
\item \emph{\Cref{eq:van-Mk}, values from the Mari\~no--Vafa side:} $G\leq 4$, $1638$ relations; $G\leq 6$ modulo $2^{61}-1$, $218+196$ relations; no failure.
\item \emph{\Cref{eq:van-Mk}, values from the Hodge side, the Mari\~no--Vafa formula not assumed:} $G\leq 5$, including $\tau=-1/3$, a zero of $\Pi_3$; $186$ relations, no failure; five kinds of control break it.
\item \emph{Target coefficient} for every part $k_0$, including repeated parts: $202/202$, $124/124$, $32/32$.
\item \emph{Mechanical run of the logic of the proof of \cref{thm:van-main}} over all terms, all $\msV$ for $G\leq 5$: $53.4$ million terms, no survivors, one target per case.
\item \emph{\Cref{thm:van-main} directly on the Hodge side:} $h=3$ ($r\leq 4$), $h=4$ ($r\leq 5$, all multisets in range), $h=5$ ($r\leq 3$), $h=6$ ($r\leq 2$), at $\tau=2/7,-5/3,-1/2,13/9$: $383$ tests, all zero; just outside the range all $228$ controls are non-zero.
\end{itemize}

\emph{Limits.} The direct comparison on the Hodge side of \cref{eq:van-real} at $g=6$ did not finish. The mechanical run at $G=6,7,8$ with $r\leq 2$ did not finish. Neither is claimed.

\subsection{Resolved conifold}
\label{sec:evi-con}

\begin{itemize}
\item Genus $2$ and $3$: the frame-B sums \cref{eq:con-frame}, exact, $\cQ$ and $\mathsf{x}$ symbolic, equal the point-frame sum at the general torus reconstructed from $97$ integer weight samples and specialised. This tests steps (d), (e) and the constants of the proof of \cref{prop:con-frame}, not (a) or (b). Steps (a)--(c) are tested independently against \cite{BS}: the degree-$1$ localisation formula in genus $1,2$, and the result on the conifold, $\sF_g-\sF_g(0)=\mathsf{P}_g\mathrm{Li}_{3-2g}(\cQ)/(\epsilon_1\epsilon_2)$ on $\TF$, in genus $2,3$.
\item Genus $4$, $27172$ decorated graphs: total minus $\mathsf{P}_4\mathrm{Li}_{-5}(\cQ)/(\epsilon_1\epsilon_2)$ is independent of $\cQ$, and equals $-(\mathsf{x}^6+1)/87091200$, the constant of the remark on $\sF_g(0)$ in \cref{sec:con}.
\item The specialisation \cref{eq:con-value} and \cref{lem:reg-H0}, in exact arithmetic from Hodge integrals, at $\mathsf{x}=2/3,3,-5/7,1$ ($h=2,3,4$; also $h=5$ at $\mathsf{x}=2/3$).
\item Exact arithmetic at three rational weight points through $\cQ^7$, $z^3$ confirms \cref{eq:con-Delta}, the ratio $\Delta_+(0)/\Delta_-(0)$, and that the solution of the row equation with constants $h_{\wts_i}$ equals the Laplace row $h_{\mathsf{nw}}\mathrm{r}^{\mathrm{pt}}_{0,i}$ and fails at $-z$. This does not test the stability argument.
\item The rank-$2$ defect $\Mdef$, built from the rows, vanishes exactly over $\bbQ$ at three samples through $z^5$, $\cQ^4$, and $z^7$, $\cQ^6$ at one of them; and in $\mathbb{F}_p$ at four samples, one through $z^9$, $\cQ^6$. Controls $\Rm^{-1}(-z)$, $\Rm(-z)$, $H\qprod\to-H\qprod$ fail; dropping $h_{\mathsf{nw}}$ gives $\Mdef=0$, as the ambiguity statement says. Localisation on $\bbP^1$ with pulled-back $\psi$ classes equals the graph sum with legs at $(g,n)=(1,1),(1,2),(2,0),(2,1),(2,2)$ through $d=3$ at two weight points; legs at $+\psi$, $\Rm^{-1}(-z)$, $\Rm(-z)$ fail.
\end{itemize}

\subsection{Limits and warnings}
\label{sec:evi-limits}

\begin{itemize}
\item \emph{Genus $2$ sees only $\Fcf(0)$ at the level of $\cO(1)$.} The polar part of $\Gsum_2$ depends on $\Fcf$ only through its value at $\xi=0$, so the control $F\to F+\xi/5$ stays bounded at $g=2$. Genus $2$ alone is not a test of $\Fcf$ beyond $\xi^0$, and genus $\geq 3$ is the first real test.
\item \emph{Sensitivity to $\Fcf$} (in $\mathbb{F}_p$, $p=2^{127}-1$). Replacing $\Fcf$ by $\Fcf+(1/9)\xi^k$ in the graphs with an edge changes the sum at valuation $k-(2g-3)$ at $g=2,3$, for $k=0,\dots,5$. So the polar part sees $\Fcf$ through $\xi^{2g-4}$ at these genera, namely $\xi^0$ at $g=2$ and $\xi^2$ at $g=3$. At $g=5$, with $356$ graphs with an edge and $9484$ coloured terms, $\xi^6$ gives poles ($3$ colourings) and $\xi^7$ gives none. With $\bw\to-\bw$ there are poles in $2926$ of $9484$ terms.
\item \emph{Rank-one sign controls do not discriminate at odd genus.} $\int_{\Mbar_g}\Hg_{\mathsf{v}}$ is even in $\mathsf{v}$ when $3g-3$ is even. So the control $\Theta(+z/\mathsf{v})$ against $\Theta(-z/\mathsf{v})$ does not fail at $g=3$. Genus-$3$ sign tests must use the total on $\TF$, where the flips do fail. Likewise, at $g=3$, $y^0$ the sign flip of $c_j$ is invisible.
\item \emph{$\Rm$ against $\Rm^*$.} $\Rm$ and $\Rm^*$ have the same $\Rm_1$. Tests that see only first order in $z$ cannot separate them. The matches at genus $\geq 2$ above do see $z^2$ and higher, but they remain evidence. Exact divisor identity $A\int_{\Mbar_{g,n+1}}\Omega(H^{n+1})=\theta\int_{\Mbar_{g,n}}\Omega(H^n)$ at $\alpha=(1,2,5)$, $\epsilon=(3,7)$: for $\Rm$ the defect is $0$ at $(g,n)=(0,4),(1,1),(2,0)$ through $y^2$ and at $(3,0)$ through $y^1$. For $\Rm^*$ the defect is non-zero already at $y^1$: $1249/2268$, $76487/17280$, $-28742463168731/35831808000$, $-406623543620925998705062409/62912004762894336000000$. $\Rm^*$ differs from $\Rm$ only in class degree $\geq 2$ of $\Omega$. $\Rm^*$ passes every test at $(g,n)=(1,1)$ and every test with total $\psi$ power $k\geq\dim\Mbar_{g,n}-1$ at $(1,2),(2,1),(2,2)$; those tests cannot separate $\Rm$ from $\Rm^*$. Below that power it fails at $(1,2)$, $k=0$ (in $3$ of $6$), $(2,1)$, $k=0,1,2$ ($2/3,3/3,3/3$), $(2,2)$, $k=0,1,2,3$ ($3/6,8/9,15/15,18/18$), first at $y^1$.
\item \emph{Scale law.} Row and edge checks at $n$ weight pairs are one test of the scale-free object plus a test of the scaling (\cref{sec:evi-rows}). Quantities on the Hodge side do depend on $\tau$.
\item \emph{Form of the coefficients.} At genus $2,3,4$ the computed $\sF_{k,h}$ have the form $\bbQ[X^{\pm1}][F]$; this is evidence only and is not proved for any $g$ in general.
\item \emph{Class-level identity.} The class-level identity $\Omega=\Rm.\TFT$ rests on \cite[Theorem~(2)]{Tel12}. The evidence pairs it with $\psi$ classes pulled back from $\Mbar_{g,n}$ at $d\leq 3$ ($g\leq 2$) and $d\leq 2$ ($g=3$, $n=1$), and with powers of $\lambda_1$ at $d\leq 2$ ($g\leq 2$). It was not paired with boundary strata or with general tautological classes.
\item \emph{Coverage.} Nothing was checked to all orders by computation. Genus: $g\leq 3$ for most items; $g=4$ with one or two pairs $\epsilon$ and points $t$. At $g\geq 5$ no full total has been computed. What exists is partial: $g=5$ per colouring for graphs with a vertex of genus $\geq 2$ at three pairs $\epsilon$; $g=6$ for about $1000$ of $5226$ such graphs at one pair $\epsilon$, a run that did not finish; the single-vertex polar part through $g=11$ at $\epsilon=(3,7)$. Degree: $d\leq 3$ ($d\leq 2$ at $g=4$; $d\leq 1$ at $g=4$ for generic $\alpha$). Order in $z$: $z\leq 3$, and $z\leq 9$ with a separate implementation, for Wick rows on $\TF$; $z\leq 5$ for rows at generic $\alpha$ against the $I$-function, extended to $z^9$, $y^6$ at one sample modulo $p$; $z^7$ for the constants $a_k$ at generic $\alpha$ through $g=4$. Order in $y$: $y^3$ at generic $\alpha$. Samples are small sets of rational weights. The refined BPS numbers are an external control, not used in any proof. Numeric regularity near $\TF$ at fixed non-zero $y$ was tested only at $g=2$. \Cref{thm:reg-main} per colouring was tested at $g=2,3$ for all graphs, with the partial coverage above at $g=5,6$.
\end{itemize}
\endgroup
